\documentclass[11pt]{article}
\usepackage[margin=0.84in]{geometry}
\usepackage{amsmath,amssymb,amsthm,mathtools}
\usepackage{graphicx}
\usepackage{microtype}
\usepackage[colorlinks=true,linkcolor=blue,urlcolor=blue,citecolor=blue]{hyperref}

\newtheorem{theorem}{Theorem}
\newtheorem{lemma}{Lemma}
\newcommand{\R}{\mathbb R}
\newcommand{\Sclass}{\mathcal S}
\newcommand{\Boxop}{\mathord\Box}

\title{The wave and Schr\"odinger Gagliardo--Nirenberg conjectures}
\author{Candidate full solution to the closing conjecture of arXiv:2201.13050}
\date{13 August 2026}

\begin{document}
\maketitle

\noindent\fbox{\parbox{0.95\linewidth}{\textbf{Status: candidate full proof,
likely valid.}  Both exponent regions conjectured by Mandel follow by
globalizing the paper's compact-piece theorem with ordinary, respectively
parabolic, Littlewood--Paley decompositions.}}

\section{Source conjecture}

Let $z=(x,t)\in\R^{d-1}\times\R$.  Mandel studies
\begin{align}
 \|u\|_q&\lesssim \|\Boxop u\|_r^{1-\kappa}\|u\|_r^\kappa,
 &\Boxop&=\partial_{tt}-\Delta_x,\label{eq:wave}\tag{W}\\
 \|v\|_q&\lesssim \|(i\partial_t-\Delta_x)v\|_r^{1-\kappa}
                    \|v\|_r^\kappa.\label{eq:sch}\tag{S}
\end{align}
The source is R.~Mandel, \emph{On Gagliardo--Nirenberg inequalities with
vanishing symbols}, arXiv:2201.13050, Anal. PDE 17 (2024), 3447--3476
\cite{Mandel}.

\begin{figure}[ht]
\centering
\includegraphics[width=0.96\linewidth]{figures/source_conjecture_crop.png}
\caption{The two conjectured exponent regions on source PDF page 29.}
\end{figure}

Precisely, for $1<r\le2\le q<\infty$ and $0<\kappa<1$, the paper conjectures
\eqref{eq:wave} under
\begin{equation}\label{eq:wave-region}
 \frac1r-\frac1q=\frac{2(1-\kappa)}d,
 \qquad
 \min\left\{\frac1r,\frac1{q'}\right\}
 \ge \frac{d-2\kappa}{2(d-1)},
\end{equation}
and \eqref{eq:sch} under
\begin{equation}\label{eq:sch-region}
 \frac1r-\frac1q=\frac{2(1-\kappa)}{d+1},
 \qquad
 \min\left\{\frac1r,\frac1{q'}\right\}
 \ge \frac{d+1-2\kappa}{2d}.
\end{equation}
The wave statement is understood for $d\ge2$, and the Schr\"odinger statement
for $d\ge2$, as in the source section.

\section{Proof intuition}

The hard oscillatory-integral work is already local in the source.  On one
Fourier annulus, the wave cone has $d-2$ nonzero principal curvatures and the
paraboloid has $d-1$.  Substituting these numbers and the vanishing orders
$(1,0)$ into the source's compact-piece theorem gives exactly
\eqref{eq:wave-region} and \eqref{eq:sch-region}.  Homogeneity makes the
constant independent of the annulus.

It remains to add the annuli without losing the multiplicative right-hand
side.  This is precisely where the conjectural restriction
$r\le2\le q$ matters.  Littlewood--Paley theory first places the output
annuli in $\ell^2$.  Sequence H\"older separates the two factors.  Finally,
$L^r$ is 2-concave for $r\le2$, while $L^q$ is 2-convex for $q\ge2$; hence
both $\ell^2$ input sums collapse to the two global norms.  The paraboloid
uses parabolic rather than isotropic annuli.  A direct one-dimensional
normal estimate handles the flat $d=2$ wave cone.

\section{Two square-function facts}

We use either the ordinary frequency dilations $\xi\mapsto\lambda\xi$ or the
parabolic dilations
\begin{equation}\label{eq:parabolic-dilation}
 \delta_\lambda(\xi',\tau)=(\lambda\xi',\lambda^2\tau).
\end{equation}
For the latter, a convenient homogeneous quasi-norm is
$\rho(\xi',\tau)=(|\xi'|^4+\tau^2)^{1/4}$.  A smooth partition of unity on
$\R^d\setminus\{0\}$ gives annular multipliers $(\Delta_j)_{j\in\mathbb Z}$.
In both dilation systems the Littlewood--Paley theorem says
\begin{equation}\label{eq:LP}
 \|f\|_p\asymp
 \left\|\left(\sum_j|\Delta_jf|^2\right)^{1/2}\right\|_p,
 \qquad 1<p<\infty.
\end{equation}
For the parabolic system this is the standard anisotropic square-function
theorem; it also follows from random signs, the anisotropic Mikhlin theorem,
and Khintchine's inequality \cite{Stein}.

We shall additionally use the elementary inequalities
\begin{align}
 \left\|\left(\sum_j|f_j|^2\right)^{1/2}\right\|_q
 &\le \left(\sum_j\|f_j\|_q^2\right)^{1/2}, &&q\ge2,
 \label{eq:2convex}\\
 \left(\sum_j\|f_j\|_r^2\right)^{1/2}
 &\le \left\|\left(\sum_j|f_j|^2\right)^{1/2}\right\|_r, &&1\le r\le2.
 \label{eq:2concave}
\end{align}
The first is the triangle inequality in $L^{q/2}$.  Raising the second to
the $r$th power reduces it to Minkowski's integral inequality in
$\ell^{2/r}$.

\section{Uniform annular estimates}

\begin{lemma}[Wave annulus]\label{lem:wave-annulus}
Suppose $d\ge3$, $0<\kappa<1$, and $1<r\le2\le q<\infty$ satisfy
\eqref{eq:wave-region}.  If
$\operatorname{supp}\widehat w\subset\{\xi:1/2\le|\xi|\le2\}$, then
\begin{equation}\label{eq:wave-annulus}
 \|w\|_q\le C\|\Boxop w\|_r^{1-\kappa}\|w\|_r^\kappa.
\end{equation}
The same constant works on every isotropic dilate of the annulus.
\end{lemma}

\begin{proof}
Away from the characteristic cone, the symbol is bounded away from zero, so
Bernstein and a smooth inverse multiplier give \eqref{eq:wave-annulus}.
Near the cone, use a finite smooth partition into compact graph pieces.  The
cone away from its vertex is smooth and has $k_0=d-2$ nonzero principal
curvatures.  Apply Theorem 4 and Remark 1(d) of \cite{Mandel} with vanishing
orders $\alpha_1=1$, $\alpha_2=0$.  Their averaged order is
$\bar\alpha=1-\kappa$, and the two local conditions become
\[
 \frac1r-\frac1q\ge\frac{2\bar\alpha}{k_0+2}
 =\frac{2(1-\kappa)}d,
 \qquad
 \min\left\{\frac1r,\frac1{q'}\right\}
 \ge\frac{k_0+2\bar\alpha}{2(k_0+1)}
 =\frac{d-2\kappa}{2(d-1)}.
\]
The non-strict inequalities are allowed because
$0<\bar\alpha<1$, $\alpha_1\ne\alpha_2$, and $0<\kappa<1$.
Finite recombination proves the estimate.  The wave symbol has isotropic
degree two, and $d(1/r-1/q)=2(1-\kappa)$, so dilation gives the same constant
on every annulus.
\end{proof}

\begin{lemma}[Schr\"odinger annulus]\label{lem:sch-annulus}
Suppose $d\ge2$ and the exponents satisfy \eqref{eq:sch-region}.  If
$\operatorname{supp}\widehat w\subset\{1/2\le\rho\le2\}$, then
\begin{equation}\label{eq:sch-annulus}
 \|w\|_q\le C\|(i\partial_t-\Delta_x)w\|_r^{1-\kappa}
                  \|w\|_r^\kappa.
\end{equation}
The same constant works on all parabolic dilates.
\end{lemma}

\begin{proof}
The characteristic paraboloid has $k_0=d-1$ nonzero principal curvatures.
The same substitution in the source theorem gives
\[
 \frac{2\bar\alpha}{k_0+2}=\frac{2(1-\kappa)}{d+1},
 \qquad
 \frac{k_0+2\bar\alpha}{2(k_0+1)}
 =\frac{d+1-2\kappa}{2d}.
\]
Away from the paraboloid use Bernstein and an inverse multiplier.  Under
\eqref{eq:parabolic-dilation}, the homogeneous dimension is
$Q=(d-1)+2=d+1$ and the Schr\"odinger symbol has degree two.  The identity
$Q(1/r-1/q)=2(1-\kappa)$ makes the constant dilation-invariant.
\end{proof}

\begin{lemma}[Flat two-dimensional wave annulus]\label{lem:flat}
Lemma \ref{lem:wave-annulus} also holds when $d=2$.
\end{lemma}

\begin{proof}
Near either characteristic line, take normal/tangent Fourier coordinates
$(\eta,\sigma)$.  On a unit annular patch, $|\sigma|$ is bounded above and
below, and the wave symbol is a nonzero constant multiple of $\eta\sigma$.
Thus $\|\partial_yw\|_r\asymp\|\Boxop w\|_r$, where $y$ is the normal physical
coordinate.  Put $a=1-\kappa=1/r-1/q$.  For every tangent coordinate $x$,
the one-dimensional Gagliardo--Nirenberg inequality gives
\[
 \|w(x,\cdot)\|_{L_y^q}
 \lesssim \|\partial_yw(x,\cdot)\|_{L_y^r}^{a}
          \|w(x,\cdot)\|_{L_y^r}^{1-a}.
\]
H\"older in $x$ and vector-valued Bernstein for the compact tangent-frequency
support bound the two resulting $L_x^q(L_y^r)$ norms by the corresponding
global $L^r$ norms.  This proves the local estimate.  Away from the lines use
the inverse-multiplier argument, and finite recombination finishes the unit
annulus.  Scale invariance gives all annuli.
\end{proof}

\section{Full solution}

\begin{theorem}\label{thm:main}
Both conjectures on source PDF page 29 are true.  More explicitly:
\begin{enumerate}
\item If $d\ge2$ and \eqref{eq:wave-region} holds, then \eqref{eq:wave} holds
for every $u\in\Sclass(\R^d)$.
\item If $d\ge2$ and \eqref{eq:sch-region} holds, then \eqref{eq:sch} holds
for every $v\in\Sclass(\R^d)$.
\end{enumerate}
\end{theorem}

\begin{proof}
First take the ordinary annular projections and write
$A_j=\|\Delta_j\Boxop u\|_r$, $B_j=\|\Delta_ju\|_r$.  Since constant
coefficient multipliers commute, Lemmas \ref{lem:wave-annulus} and
\ref{lem:flat}, \eqref{eq:LP}, and \eqref{eq:2convex} give
\[
 \|u\|_q
 \lesssim\left(\sum_j A_j^{2(1-\kappa)}B_j^{2\kappa}\right)^{1/2}.
\]
H\"older in the sequence variable, followed by \eqref{eq:2concave} and
\eqref{eq:LP}, yields
\begin{align*}
 \|u\|_q
 &\lesssim
 \left(\sum_jA_j^2\right)^{(1-\kappa)/2}
 \left(\sum_jB_j^2\right)^{\kappa/2}\\
 &\lesssim \|\Boxop u\|_r^{1-\kappa}\|u\|_r^\kappa.
\end{align*}
This is \eqref{eq:wave}.

For \eqref{eq:sch}, repeat the identical computation with the parabolic
annular projections and Lemma \ref{lem:sch-annulus}.  The anisotropic version
of \eqref{eq:LP} supplies both square-function bounds.  This proves the
Schr\"odinger assertion.
\end{proof}

\section{Verification, scope, and novelty}

The proof uses the source's own proved compact-piece theorem, plus standard
square-function theory.  A deterministic check in the packet verifies both
curvature substitutions, both homogeneous dimensions, and the sequence
inequalities on randomized finite data.  The LaTeX packet was rendered and
visually inspected page by page.

The theorem covers the full stated ranges, including equality in the
curvature thresholds.  It does not address $r=1$, $q=\infty$, operators with
general unbounded characteristic sets, optimal constants, or extremizers.

Bounded novelty searches on 13 August 2026 covered the run indexes, the exact
title and conjecture language, wave/Schr\"odinger Gagliardo--Nirenberg phrases,
the displayed exponent forms, and later works citing the source.  No later
primary answer was found.  Novelty confidence is moderate: the globalization
is short and may be known informally.

\textbf{Human-review recommendation.}  Verify that Theorem 4 plus Remark 1(d)
of \cite{Mandel} applies uniformly to the finitely many unit-annulus cone and
paraboloid patches, and check the anisotropic Littlewood--Paley normalization.
The global summation is elementary once those annular estimates are accepted.

\begin{thebibliography}{9}
\bibitem{Mandel}
R.~Mandel,
\emph{On Gagliardo--Nirenberg inequalities with vanishing symbols},
Anal. PDE \textbf{17} (2024), 3447--3476;
\url{https://arxiv.org/abs/2201.13050}.

\bibitem{Stein}
E.~M. Stein,
\emph{Harmonic Analysis: Real-Variable Methods, Orthogonality, and
Oscillatory Integrals}, Princeton Mathematical Series 43,
Princeton University Press, 1993.
\end{thebibliography}

\end{document}
